\subsection{The entire principal part and the cost of lower coefficients}
\label{sec:fulljet}
The highest coefficient is only one part of a logarithmic soft
expansion. In a fixed reference-energy convention, write
$\beta_k=k!u_k$, $0\le k\le m$, for all coefficients of its Mellin
principal part. Lower coefficients can depend on the higher ones
under a change of convention, as Section~\ref{sec:units} explains.
We now give simultaneous estimators and componentwise optimal rates
for this full principal part. Holomorphic Laurent coefficients are
not recovery targets here.

At the same scales as Theorem~\ref{thm:vectorlogs}, put
$L=-\log h$ and define the Lagrange polynomials
\[
 \ell_j(t)=\prod_{\substack{0\le r\le m\\r\ne j}}
 \frac{t-L-rd}{d(j-r)},\qquad 0\le j\le m.
\]
For coefficient extraction denoted by $[t^n]$, set
\begin{equation}\label{eq:fullweights}
 W_{k,j}=\sum_{r=0}^{\lfloor(m-k)/2\rfloor}
 \frac{(-\sigma^2/2)^r}{r!}(k+2r)!
 [t^{k+2r}]\ell_j(t),\qquad
 \widehat\beta_k=\sum_{j=0}^m W_{k,j}E_{he^{-jd},\sigma}(Y).
\end{equation}
These are explicit finite real weights. No unknown soft coefficient
enters their construction. For $n\ge0$ define the positive polynomial
\[
 H_n(A;\sigma)=\sum_{r=0}^{\lfloor n/2\rfloor}
 \frac{(\sigma^2/2)^r A^{n-2r}}{r!(n-2r)!}.
\]

\begin{theorem}[Full principal-part recovery]\label{thm:fulljet}
Under \eqref{eq:vectorclass}, let $m\ge1$ and fix $d,\sigma>0$.
For every $k\le m$, the estimator \eqref{eq:fullweights} satisfies
\begin{align}
 \|\widehat\beta_k-\beta_k\|_X
 &\le \frac{H_{m-k}(|L|+md;\sigma)}{d^m}\,
 \mathcal E(h),\label{eq:fullbound}\\
 \mathcal E(h)&=
 e^{\alpha^2\sigma^2/2}Mh^\alpha(1+e^{-\alpha d})^m
 +\frac{e^{a^2\sigma^2/2}}{\sqrt{2\pi}\sigma}
 \eta h^{-a}(1+e^{ad})^m.\nonumber
\end{align}
Let $\mathcal R_{m,k}(\eta)$ be the minimax norm error for the
individual target $\beta_k$. As $\eta/M\downarrow0$, with all other
parameters fixed,
\begin{equation}\label{eq:fullrate}
 \mathcal R_{m,k}(\eta)\asymp
 M^{a/(a+\alpha)}\eta^{\alpha/(a+\alpha)}
 [\log(M/\eta)]^{m-k}.
\end{equation}
Here $\asymp$ means upper and lower positive constants independent of
$M,\eta$ in the stated small-noise regime. It does not assert a sharp
minimax constant or a bound uniform in $m$.
\end{theorem}
\begin{proof}
The polynomial component of each exact-data Gaussian average is
$S(t)=\mathbb E P(t-\sigma Z)$, where $P(t)=\sum u_kt^k$.
On degree-at-most-$m$ polynomials, Gaussian averaging is the finite
heat operator $S=e^{(\sigma^2/2)\partial_t^2}P$.
Lagrange interpolation reconstructs $S$ from $S(L+jd)$, and the
finite inverse operator $e^{-(\sigma^2/2)\partial_t^2}$ recovers $P$.
Applying $k![t^k]$ gives exactly \eqref{eq:fullweights}.
This argument uses the polynomial component of a convergent Gaussian
average; it does not assume that a global polynomial alone has an
absolutely convergent Mellin transform.

Put $A=|L|+md$. The denominator of $\ell_j$ has absolute value
$d^m j!(m-j)!$, and every root of its numerator has absolute value
at most $A$. Hence, for $0\le n\le m$,
\[
 |[t^n]\ell_j(t)|\le
 \frac{\binom mn A^{m-n}}{d^m j!(m-j)!}.
\]
Substitution in \eqref{eq:fullweights} gives
\[
 |W_{k,j}|\le\binom mj d^{-m}H_{m-k}(A;\sigma).
\]
Each observed Gaussian average differs from its polynomial component
by at most
\[
 e^{\alpha^2\sigma^2/2}Mh^\alpha e^{-\alpha jd}
 +\frac{e^{a^2\sigma^2/2}}{\sqrt{2\pi}\sigma}
 \eta h^{-a}e^{ajd}.
\]
Summation proves \eqref{eq:fullbound}.

For the upper rate choose
$h=(\eta/M)^{1/(a+\alpha)}$ and
$L=\log(M/\eta)/(a+\alpha)$.
Both terms in $\mathcal E(h)$ are a fixed constant times
$S_\eta=M^{a/(a+\alpha)}\eta^{\alpha/(a+\alpha)}$.
Since $H_{m-k}(|L|+md;\sigma)=O((1+|L|)^{m-k})$, the upper
bound in \eqref{eq:fullrate} follows.

For the lower rate use the same cutoff pair as
Theorem~\ref{thm:loglower}. Set
\[
 C_m=(m/(\alpha e))^m,\quad
 K_0=\frac{2a^{m+1}C_m}{m!},\quad
 h_0=(K_0\eta/M)^{1/(a+\alpha)},\quad L_0=-\log h_0.
\]
The coefficient of $t^k$ in
$P(t)=c(t-L_0)^m v$, with $c=Mh_0^\alpha/C_m$, is
$c\binom mk(-L_0)^{m-k}v$.
The common midpoint observation has noise at most $\eta$, so
\begin{equation}\label{eq:fulllower}
 \mathcal R_{m,k}(\eta)\ge
 \frac{1}{2(m-k)!}
 \left(\frac{m!}{C_m}\right)^{a/(a+\alpha)}
 (2a^{m+1})^{\alpha/(a+\alpha)}S_\eta |L_0|^{m-k}.
\end{equation}
Since $L_0=[\log(M/\eta)-\log K_0]/(a+\alpha)$, this proves
the matching small-noise lower rate. The zero-noise lower bound
is trivially zero.
\end{proof}

For example, at degree $m=2$, setting
$S_\eta=M^{a/(a+\alpha)}\eta^{\alpha/(a+\alpha)}$ gives
\begin{center}
\begin{tabular}{ll}
\toprule
Target coefficient & Componentwise minimax order as $\eta/M\downarrow0$\\
\midrule
Triple pole $\beta_2$ & $S_\eta$\\
Double pole $\beta_1$ & $S_\eta\log(M/\eta)$\\
Ordinary residue $\beta_0$ & $S_\eta[\log(M/\eta)]^2$\\
\bottomrule
\end{tabular}
\end{center}
This loss concerns unknown lower and higher coefficients in the stated
class, at fixed convention. Knowing some nuisance coefficients can
improve a different recovery problem. A bounded angular pairing
inherits each upper estimate with its dual norm and the corresponding
lower rate when the functional is nonzero. For a continuum frequency
band, the same weight bound applies after replacing the noise term
$\eta$ in $\mathcal E(h)$ by
$\eta\operatorname{erf}(\sigma\Omega/\sqrt2)
+Q\operatorname{erfc}(\sigma\Omega/\sqrt2)$.

Moving the interpolation window does not remove the lower bound for a
coefficient at a prescribed reference scale. Choosing a new reference
can instead change that lower coefficient, hence change the recovery
target. This is separate from a passive change of energy unit with a
co-transforming reference, which Section~\ref{sec:units} treats.
